%The main goal of the Compact Muon Solenoid (CMS) experiment, one of the detectors of the Large Hadron Collider (LHC) at CERN , is to explore particle physics at the TeV energy scale. The electromagnetic calorimeter (ECAL), one of the sub-detectors of CMS, measures the energy of electrons and photons produced in LHC collisions \cite{cmstp}. 
After Run-3 of the LHC a significant upgrade is planned for the HL-LHC \cite{Apollinari:2015wtw} period of operation. Expected instantaneous luminosities in the HL-LHC era will be a factor of five times larger than the current LHC nominal conditions. While these more intense conditions will enable the LHC experiments to accumulate more than an order of magnitude larger dataset than all of the preceding LHC runs, it will also exacerbate the experimental challenge of the multiple pile-up interactions appearing in each event. To mitigate degradation of LHC event reconstruction performance resulting from the potentially hundreds of pile-up (PU) interactions expected in an HL-LHC event, both CMS \cite{Butler:2019rpu} and ATLAS \cite{Morgenstern:2019xto} are developing methods to incorporate precision timing measurements into the detectors, significantly expanding the capabilities of the experiments. The inclusion of precision timing measurements through a combination of new sub-detectors (MTD, HGCAL) and dedicated upgrades to existing detectors, such as the CMS electromagnetic calorimeter (ECAL). This will also facilitate the use of space and time ( 4D ) in the reconstruction of events. 4D reconstruction methods will provide new tools for both pile-up suppression and a range of new experimental techniques to probing physics signatures and long-lived particles. The ECAL has set a goal of upgrading its systems to achieve $\sim$30 ps time resolution in the HL-LHC era \cite{Bruneliere:2006ra}.

%In this paper we study the timing reconstruction of the ECAL in order to improve the timing resolution of the ECAL.
The ECAL is a hermetic homogeneous calorimeter made of 75 848 lead tungstate (PbWO$^{4}$) scintillating crystals arranged in a barrel (EB) and two endcap (EE) sections \cite{cmstdr}.  Lead tungstate has a fast scintillation response and is resistant to radiation.  These crystals have a high density of 8.3 g/cm$^{3}$, a short radiation length of X$_{0}$ = 0.89 cm, and a small Moli\`ere radius of R$_{M}$ = 2.0 cm, allowing for a highly granular, compact detector. Each individual crystal is a truncated pyramid, with a lateral size comparable to R$_{M}$ and a length of 25.8 X$_{0}$ for the barrel, and 24.7 X$_{0}$ for the endcaps. When an electromagnetic particle interacts with a crystal, and thereby deposits energy within the crystal, a burst of scintillating light is emitted. The scintillation decay time of the crystals is comparable to the LHC bunch crossing interval of 25 ns in which 80\% of the light is emitted. The intensity of the scintillating light bursts are measured with avalanche photodiodes (APD) in the barrel and vacuum phototriodes (VPT) in the endcaps \cite{Bayatian:2006nff}.

%In addition to the energy measurement, the combination of the scintillation timescale of PbWO$^{4}$, the electronic pulse shaping, and the sampling rate allow excellent time resolution to be obtained with the ECAL. 
%While the primary purpose of the electromagnetic calorimeter is to measure energy deposited in the calorimeter by incident photons and electrons, the combination of the scintillation timescale of PbWO$^{4}$ crystals used in the CMS ECAL, the electronic pulse shaping, and the sampling rate allow excellent time resolution to be obtained with the CMS ECAL.

While the primary purpose of the ECAL is to measure energy deposited in the calorimeter by incident photons and electrons, the combination of the scintillation timescale of PbWO$^{4}$ crystals, the electronic pulse shaping, and the sampling rate, allow excellent time resolution to be obtained. The time resolution achievable in the ECAL directly impacts the energy resolution and plays a significant role in many important physics measurements. Currently, time measurements allow the rejection of backgrounds with a broad time distribution such as those arising from cosmic rays, beam halo muons, electronic noise, secondary interactions within the material of the support structures and electronics of the detectors, and out-of-time proton-proton interactions \cite{Bruneliere:2006ra}. The timing resolution of the ECAL also directly impacts the sensitivity of many types of searches, such as those involving long lived particle (LLP) decays. When a particle does not decay promptly, the arrival time at the ECAL of particles resulting from the decay of this particle can be measurably later than those coming from the decay of short lived particles. 
%Time resolutions of $\sim$100 ps in a single crystal have been measured in Run 2 operation \cite{Paganini:2004uw}.

%The current time reconstruction is sensitive to the presence of pulses arising from particles produced in proton-proton collisions that occurred before or after the current collision.  

%The time reconstruction is sensitive to the presence of scintillating light pulses in a crystal that arise from particles produced in collisions other then the collision that produced the particle interacts that are being reconstructed within ECAL. 
The time reconstruction is sensitive to the scintillating light pulses in a crystal that arise from particles produced in collisions other than those of the particles produced in the particular collision that are to be reconstructed. Such extraneous collisions present in a reconstruction are called pile-up (PU) collisions. When these PU collisions originate in the bunches of colliding protons proceeding or following the bunch of colliding protons that produced the reconstructed interactions, this is referred to as out-of-time pile-up (OOT-PU). As the LHC moves to higher luminosity in the HL-LHC era, the effects of OOT-PU in events become a significant issue in many areas.  This will not only directly degrade the time resolution, but since the time measurement is used to reject backgrounds arising from OOT-PU, the resolution of many physics measurements can be degraded. A significant step to prepare for operations in high OOT-PU environments will be to study methods to mitigate the effects of OOT-PU on the ECAL timing reconstruction.  A timing reconstruction method that is OOT-PU aware will result in an improved time resolution that will not only improve the precision of many current physics measurements, it will also provide a foundation to meet the goal of achieving a $\sim$30 ps time resolution in the ECAL. Such a timing method will support the incorporation of precision timing in the HL-LHC era, and will also allow us to begin to explore in Run III some of the possibilities of 4D reconstruction methods facilitate by precision timing measurements in the HL-LHC era.

%Improved time resolution also allows us to better resolve vertexes in an event by looking at the possible intersection of reconstructed particle paths based on their possible path length. The resolution for the location of such vertexes comes directly from the resolution of the path lengths, which are determined from the differences in arrival time. In the case where the reconstructed particles are photons, the timing information from the ECAL is the only way to reconstruct possible secondary vertexes.

%In this note we will present the methodology which was used to determine the time resolution of the ECAL that will be used to study the effectiveness of proposed new timing algorithms to mitigated the effect of OOT PU.  

In this paper a we will discuss time reconstructions methods in Section 2. First, in Section 2.1 we will review the preceding time reconstruction, the ratio method, and the challenges faced with this reconstruction. We will then present, in Section 2.2 in the details of new the time reconstruction, the Cross Correlation method. In this section we will also discuss the nature of, and the setting of the parameters used to reject, anomalous reconstructed hits, called "spikes", in Section~\ref{sec:resspkdata}. We will also discuss the sensitivity of the new time reconstruction on out-of-time amplitudes derived in the energy reconstructed in Section~\ref{sec:recobias}.  Finally, the performance of the Cross Correlation time reconstruction in data is presented in Section 3.  Finally, this paper is summarized in Section 4.

%the methodology which was used to determine time resolution of the ECAL. With this background we will then more fully discuss the challenges that the current timing reconstruction method faces as the operation of the LHC has increased in luminosity, and then we will present the implementation of a new timing algorithm to meet these challenges.  We then present the implementation a MC truth time using the pCalo hit information in order to determine the accuracy of the new timing algorithm. Finally, we will review and contrast the performance of this timing algorithm and the ratio method, which is the currently utilized time reconstruction algorithm.


